\section{Related Work}
\label{sec:relatedwork}

\subsection{FAR benchmarks and the two-axis evaluation status quo}
FAR-Trans~\cite{sanzcruzado:fartrans} is the first public FAR benchmark with both real retail-investor transactions and asset pricing, and it establishes the evaluation standard for the field: 11 baselines are reported on nDCG@10 for next-purchase ranking quality and ROI@10 for the realised return. The subsequent study by \citeauthor{sanzcruzado:investors}~[\citeyear{sanzcruzado:investors}] shows that these two axes are not aligned on FAR-Trans: transaction-based and profitability-based metrics are negatively correlated, and models that win one lose the other. However, neither axis measures whether a recommendation aligns with the customer's declared MiFID II risk band, which is the gap PC@$k$ addresses.

\subsection{Profile-aware FAR architectures}
\citeauthor{ghiye:rolling-forward}~[\citeyear{ghiye:rolling-forward}] extend LightGCN with a causal graph convolution and a sliding-window training scheme for credit-bond recommendation, capturing temporal drift in user interests via embedding updates restricted to recent interactions. Their result supports the broader point that LightGCN can be modified to carry domain-specific structure beyond plain Bayesian Personalised Ranking (BPR), and motivates the choice of LightGCN as the backbone for our stratified, profile-coherent extension.

\subsection{Risk-aware FAR via post-hoc re-ranking}
\citeauthor{sakurai:risk-aware}~[\citeyear{sakurai:risk-aware}] propose a risk-aware utility re-ranking (RURA) procedure on FAR-Trans: a mean-variance utility conditioned on the customer's MiFID risk tolerance is applied to a base ranker's top list, lifting ROI@10 while preserving nDCG@10. RURA optimises a risk-tolerance utility but does not quantify how aligned the resulting list is with the declared band in terms of ordinal distance. PC@$k$ is complementary: it is a metric for that alignment, and is computable on any ranker (with or without re-ranking).

All in all, existing FAR work optimises profitability or ranking, with risk-tolerance entering only as a post-hoc utility weight. No prior FAR evaluation metric operationalises MiFID II suitability as an ordinal coherence score over top-$k$ lists. This paper closes that gap with PC@$k$ and a profile-coherent training-time objective.
